% !PRINTED 318
\weberpage{318}
%<<CONT
$T$ is also intransitive. From this it follows, as shown above, that $TQ$
is imprimitive.

As the group for primitive metacyclic equations of degree 8
there remain, therefore, only the two types
\[
T_7 Q, \quad T_{21} Q
\]
In the first case the quantities $v$ occurring in the formula [\S.\,82, (3)]
are the roots of a cyclic equation of degree 7
with the group $(z, z+b)$, in the second of a metacyclic one with the
group $(z, az+b)$ (Bd. I, \S.\,180), in which $z, a, b$ are taken according to the modulus 7
and $a$ runs through the values 1, 2, 4.

In order to obtain an assignment of the $v_i$, as they occur in the formula (3), \S.\,82,
to the $v_z$, in which $z$ is the index occurring in the groups $(z, z+b)$ or $(z, az+b)$,
we can start from an arbitrary
substitution of the seventh order, say from
\[
\tau = \begin{pmatrix} 1, & 0, & 1 \\ 1, & 0, & 0 \\ 0, & 1, & 0 \end{pmatrix}
\quad [\S.\,79, (2)].
\]

If we then take $v_{1,0,0}$ for $v_0$, one finds by repeated
application of $\tau$ to $1, 0, 0$ the following arrangement:
\[
(1,0,0)=0, \quad (1,1,0)=1, \quad (1,1,1)=2, \quad (0,1,1)=3, \quad (1,0,1)=4,
\]
\[
(0,1,0)=5, \quad (0,0,1)=6,
\]
and by \S.\,82, (2) one has therefore to replace
\[
v_1, v_2, v_3, v_4, v_5, v_6, v_7
\]
by
\[
v_0, v_5, v_6, v_3, v_4, v_1, v_2
\]
By this there results from \S.\,82, (3):
\[
8X = A \sum_{0,\,6}^{z} \sqrt{v_z}\,\sqrt{v_{z+1}}\,\sqrt{v_{z+2}}\,\sqrt{v_{z+4}} + B,
\]
in which $A, B$ are rational quantities, and $v_z$ are the roots of a cyclic or
of a metacyclic equation of the group $(z, az+b)$ of degree 7.

Conversely, it is also not difficult to prove that the eight quantities
derived from this by change of the signs are always the roots
of a primitive metacyclic equation
%CONT>>
